% Pursuing Stacks (cotes 134-2 to 134-8): archive cote 134-4, site page 13 (PDF page 14), typescript page 101.
% typescript. Draft compared: ps3.tex:635-704.
% AI-transcribed by Claude Opus 5.5 (claude-opus-5-5, adaptive thinking, effort high), Message Batches API, batch ps1 (msgbatch_01Rv5hCaNTtUt4JHx1yqoHPL), 2026-11-12; image review recorded in records/CHECK_DECISIONS_134_3_4.tsv.
\scanpage{134-4/13}
\mspage{101}

$M\hand{\longrightarrow} W^{-1}M=H_M$ factors into
\[ M\hand{\longrightarrow} M_{\underline{I}}\hand{\longrightarrow} H_M = W^{-1}M \quad , \]
and from b) it follows that if $f,g : X\hand{\rightrightarrows} Y$ are in the same \underline{I}-class, then
f is in W iff g is, hence by passage to quotient a subset
\[ W_{\underline{I}}\hand{\subset} \mathrm F\over{\mathrm l}{l}(M_{\underline{I}}) \]
\over[typed]{i}{o}\over[typed]{n}{f} the set of arrows in $M_{\underline{I}}$ , namely a notion of \xout{qua} weak equivalence in
$M_{\underline{I}}$ . It is evident from the universal property of $H_M$ that the canonical
functor $M_{\underline{I}}\hand{\to} H_M$ induces an isomorphism of categories
\[ H_{M_{\underline{I}}} \overset{\mathrm{def}}{=} W_{\underline{I}}^{-1}M_{\underline{I}} \hand{\xrightarrow{\ \sim\ }} H_M = W^{-1}M \quad . \]

\margin{\diagram{\begin{tikzpicture}[x=1ex,y=1ex]\draw (0,0) .. controls (.5,-.1) and (.55,.4) .. (.65,1.3)--(.8,2.3) .. controls (1.3,2.2) and (1.8,2.2) .. (2.1,2.35);\end{tikzpicture}}}It's hard at this point not to expect that $W_{\underline{I}}$ should satisfy the same mild
saturation conditions as W, so let's look into this in the \over[typed]{S}{s}tri\over[typed]{\unsure{l}}{d}\over{e}{e} (even
though I have not had any use of this so far). Condition b) of saturation\hand{,}
namely that if f,g are composable and two among f,g,gf are in W so is the
third, carries over trivially. Condition a), namely the tautological
looking condition that W should contain all isomorphisms, makes already
a problem however. It is OK though if W satisfies the following \inserted[typed]{saturation} condition\hand{,}
which is a strengthening of condition c) of page 59 :

c') Let $f\colon X\hand{\longrightarrow} Y$ and $g\colon Y\hand{\longrightarrow} X$ such that $gf\hand{\in} W$ and $fg\hand{\in} W$, then
$f,g\hand{\in} W$ .

This condition c') carries over to $M_{\underline{I}}$ trivially. This suggests to intro-%
duce a strengthening of the ``mild saturation conditions'', which \xout{we} I in-%
tend henceforth to call by the name of ``\uline{saturation}'', reserving the term
of ``\uline{strong saturation}'' to what I have previously referred to \unsure{accasionally|occasionally}
as ``saturation'' - namely the \inserted{still more} exacting condition that W consists of \uline{\over{all}{all}}
arrows made invertible by the localization functor $M\hand{\longrightarrow} H_M=W^{-1}M$, or
equivalently, by some functor $M\hand{\to} H$ . \inserted{Thus} We'll say W is \uline{saturated} if it
satisfies the following :

a') For any X in M, $\mathrm{id}_X\hand{\in} W$ .

b') Same as b) before: if two among f,g,gf are in W, so is the thi\hand{rd}

c') If $f\colon X\hand{\to} Y$ and $g\colon Y\hand{\to} X$ are such that $gf,\over{gf}{gf}\hand{\in} W$, then $f,g\hand{\in} W$ .

Each of these conditions carries over from W to a $W_{\underline{I}}$ trivially.

I can't help either having a look at the most evident exactness pro-%
perties of the canonical functor $M\hand{\to} M_{\underline{I}}$ \unsure{?} Thus one immediately sees
that for two maps
\[ f,g : X\hand{\rightrightarrows} Y=Y_1\mathrm{x}Y_2 \quad \unsure{\mathrm n}, \]
with components $f_i$ , $g_i$ ($i\hand{\in}\hand{\{1,2\}}$), f and g are I-homotopic \inserted{in the wide sense} iff so are
$f_i$ and $g_i$ (for $i\hand{\in}\hand{\{1,2\}}$). The analoguous statement is valid for \xout{any} maps
into any product object $Y=\hand{\prod} Y_i$ \deleted{- not necessarily} on a finite set of

\margin{13}
